% !TeX root = ../thesisv2.tex

\begin{abstract}
高分辨率遥感影像道路提取是数字地图生产、智慧交通系统建设、灾害应急制图、国土空间治理和自动驾驶先验地图构建中的基础技术环节。与一般区域目标分割相比，遥感道路提取同时具有细长结构明显、尺度变化显著、背景干扰复杂和整图连通性要求较高等特点。道路目标常与建筑阴影、停车场纹理、裸地边缘和植被遮挡区域在局部外观上产生混淆，模型既需要保持浅层边缘细节，又需要建立足够强的多尺度上下文表达能力。因此，围绕稳定基线模型开展结构设计、统一训练和整图评估，对于形成可靠的道路提取结论具有直接意义。

本文基于当前本地全部可复核材料，重建并完善了一套完整的深度学习实验工程，形成了覆盖数据整理、离线 patch 生成、模型训练、整图推理、定量评估、结果可视化和论文写作支撑的统一研究流程。工程层面，本文将 Massachusetts Roads Dataset 统一整理为训练集、验证集和测试集目录结构，将原始遥感图像补边至 $1536 \times 1536$，在训练阶段切分为 $512 \times 512$ patch，并在验证与测试阶段采用整图切块预测、空间位置拼接恢复和整图层面的统一指标统计。通过离线 patch 数据集构建、多进程数据加载和日志同步方式优化，基线模型单 epoch 训练耗时由早期在实现中的 409.3 s 降低至 113.5 s，工程训练吞吐获得明显改善，为后续多模型完整训练和对比分析提供了稳定支撑。

方法层面，本文以基于 Batch Normalization 和双线性上采样的 U-Net 作为统一 baseline，围绕跳跃连接特征增强和瓶颈层上下文增强两条主线开展结构设计。主线模型 DLGU-Net 在 U-Net 四级跳跃连接位置引入轻量级双线性注意力模块 DLAM，通过通道分支和空间分支联合完成浅层特征重标定，以增强道路相关响应并抑制复杂背景干扰。与此同时，本文围绕 baseline 构造了空洞卷积增强基线，在瓶颈层引入并行膨胀率为 1、2、4 和 8 的多分支上下文模块，以提高模型对主干道路与细小支路的多尺度感知能力。为形成更完整的结果矩阵，本文还统一实现并训练了 Attention U-Net、UNet++、DDU-Net、ResNet34 U-Net、Vanilla U-Net 和 Residual Vanilla U-Net 等模型，并全部纳入相同整图评估协议下进行比较。

评价层面，本文全部核心结论均建立在整图验证与整图测试结果基础上，采用交并比 IoU、F1 分数、Precision、Recall 和 Accuracy 作为主要评价指标。实验结果表明，统一 baseline U-Net 在测试集上达到 IoU 0.645964、F1 0.782378，构成了参数规模适中、综合性能稳定的统一参考模型。主线模型 DLGU-Net 在约 8.30M 参数规模下达到 IoU 0.641934、F1 0.779244，体现出轻量级跳跃连接注意力增强在道路提取任务中的有效建模能力。围绕 baseline 的直接改进中，空洞卷积增强基线在测试集上达到 IoU 0.646839、F1 0.782933，相对 baseline 获得小幅提升，并成为 baseline 改进分支中的最佳测试结果。作为扩展探索分支，Vanilla U-Net 获得了全部模型中的最高测试集 IoU 0.647187 和 F1 0.783246，进一步展示了原始风格卷积块与可学习上采样在当前任务条件下的性能潜力。

本文工作的主要贡献体现在三个方面。第一，完成了原始实验工程的系统重建，并通过离线 patch 数据集和多进程数据加载显著降低了训练阶段的数据瓶颈，建立了适合多模型统一评估的实验平台。第二，围绕 U-Net 主干形成了以 DLGU-Net 为主线模型、以空洞卷积增强基线为直接改进分支、以 Vanilla U-Net 等为扩展探索分支的完整实验矩阵，保证了模型分析具有足够广度。第三，将全部结论统一建立在整图验证和整图测试协议之上，形成了数据导向、结构清晰、可复核的模型比较体系。研究结果表明，在当前数据规模和任务特征下，以稳定 baseline 为锚点开展轻量、直接、可控的结构增强，是获得持续性能收益的有效路径。

\xtusetup{
  keywords = {遥感影像, 道路提取, 语义分割, U-Net, 注意力机制, 空洞卷积},
}
\end{abstract}

\begin{abstract*}
Road extraction from high-resolution remote sensing imagery is a fundamental task in digital map production, intelligent transportation systems, disaster-response mapping, land-space governance, and prior-map construction for autonomous driving. Compared with ordinary region segmentation, remote-sensing road extraction involves elongated structures, strong scale variation, complex background interference, and high requirements on full-image connectivity. Roads are frequently confused with building shadows, parking-lot textures, bare-soil boundaries, and vegetation occlusions. Therefore, a practical model must preserve shallow boundary details while simultaneously building strong multi-scale contextual representation. Under this problem setting, structural exploration around a stable baseline together with unified full-image evaluation is essential for obtaining reliable conclusions.

Based on all currently available local materials and fully reproducible experimental records, this thesis reconstructs and consolidates a complete deep learning workflow for road extraction. The workflow covers dataset organization, offline patch generation, model training, full-image inference, quantitative evaluation, visualization, and thesis-oriented result management. At the engineering level, the Massachusetts Roads Dataset is reorganized into a unified train/validation/test structure. Each raw image is padded to $1536 \times 1536$, split into $512 \times 512$ patches during training, and restored by tiled inference and reconstruction during validation and testing. Through offline patch generation, multi-worker data loading, and logging optimization, the single-epoch training time of the baseline model is reduced from 409.3 s in the early implementation to 113.5 s in the optimized pipeline, which makes large-scale multi-model comparison feasible.

At the methodological level, a U-Net with Batch Normalization and bilinear upsampling is adopted as the unified baseline. Two structural routes are then investigated. The first route focuses on skip-feature enhancement. The main model DLGU-Net introduces a lightweight dual-linear attention module, DLAM, into the four skip connections of U-Net. DLAM jointly performs channel recalibration and spatial-direction enhancement so that road-relevant responses can be strengthened before being fused into the decoder. The second route focuses on bottleneck-context enhancement. A dilated-convolution-enhanced baseline is designed by inserting four parallel atrous branches with dilation rates 1, 2, 4, and 8 into the bottleneck, aiming to improve the perception of major roads and narrow branches at multiple scales. In addition, Attention U-Net, UNet++, DDU-Net, ResNet34 U-Net, Vanilla U-Net, and Residual Vanilla U-Net are also implemented and trained under the same protocol to form a complete comparison matrix.

All final conclusions in this thesis are derived from unified full-image validation and testing. The evaluation metrics include IoU, F1-score, Precision, Recall, and Accuracy. Experimental results show that the unified baseline U-Net achieves a test IoU of 0.645964 and an F1-score of 0.782378, serving as a stable and moderate-scale reference model. The main model DLGU-Net reaches a test IoU of 0.641934 and an F1-score of 0.779244 with about 8.30M parameters, demonstrating the effectiveness of lightweight skip-feature attention in remote-sensing road extraction. Among the direct baseline-oriented improvements, the dilated-convolution-enhanced baseline achieves a test IoU of 0.646839 and an F1-score of 0.782933, yielding a slight improvement over the unified baseline. In the extended exploration branch, Vanilla U-Net attains the best overall test IoU of 0.647187 and the best F1-score of 0.783246, showing the performance potential of original-style convolution blocks and learnable upsampling under the current task setting.

The contributions of this thesis are threefold. First, the original project is systematically rebuilt into a reusable experimental platform, and the data bottleneck is significantly reduced through offline patches and multi-worker loading. Second, a structured experimental matrix is established around the U-Net backbone, with DLGU-Net as the main model, the dilated baseline as the direct enhancement branch, and Vanilla U-Net together with other variants as the extended exploration branch. Third, all conclusions are consistently derived from unified full-image validation and testing, resulting in a data-driven, structurally clear, and reproducible comparison framework. The results indicate that, under the current data condition and task characteristics, lightweight, direct, and controllable enhancement around a stable baseline is an effective strategy for obtaining sustained performance gains in remote-sensing road extraction.

\xtusetup{
  keywords* = {remote sensing imagery, road extraction, semantic segmentation, U-Net, attention mechanism, dilated convolution},
}
\end{abstract*}
